\documentclass[9pt,a5paper]{extarticle}
\usepackage[left=0.6in,right=0.6in,top=0.6in,bottom=0.6in]{geometry}
\usepackage{amsmath,amssymb,mathtools}
\usepackage{graphicx}
\usepackage{float}
\usepackage{caption}
\usepackage{listings}
\usepackage{xepersian}
\settextfont[Path=../fonts/,Scale=1,BoldFont={XB NiloofarBd.ttf},ItalicFont={XB NiloofarIt.ttf},BoldItalicFont={XB NiloofarBdIt.ttf}]{XB Niloofar.ttf}
\linespread{1.3}
\captionsetup{font=small}
\graphicspath{{figures/}}
\emergencystretch=2em
\newcommand{\shot}[2]{\begin{figure}[H]\centering\includegraphics[width=0.95\linewidth]{#1}\caption{#2}\end{figure}}
\newcommand{\placeholder}{\textit{این بخش پس از انجام آزمایش تکمیل می‌شود.}}

\begin{document}
	\centerline{\Large\bf گزارش آزمایش ۹: پیاده‌سازی حافظه‌ی \lr{TCAM}}
	\smallskip
	\begin{center}
		\begin{tabular}{cc}
			پوریا عرفانزاده & 401110918\\
			فاطمه ابراهیمی & 401110931
		\end{tabular}
	\end{center}
	\vskip30pt plus 5pt

	\section*{۹.۱ مقدمه و اهداف}
	در این آزمایش، یک حافظه‌ی آدرس‌پذیر با محتوا از نوع سه‌حالته طراحی می‌شود. پیمانه‌ی \lr{TCAM} شامل ۱۶ ثبات ۱۶ بیتی است و علاوه بر صفر و یک، حالت بی‌تفاوت \lr{X} را هنگام تطبیق پشتیبانی می‌کند.

	\section*{۹.۲ نمایش داده و منطق تطبیق}
	\placeholder

	\section*{۹.۳ طراحی و کد \lr{Verilog}}
	\placeholder

	\section*{۹.۴ تست‌بنچ و شبیه‌سازی}
	\placeholder

	\section*{۹.۵ سنتز و پیاده‌سازی روی برد}
	\placeholder

	\section*{نتیجه‌گیری}
	\placeholder
\end{document}
